\section{Empirical Contribution}\label{sec:empirical}

\subsection{Falsification of BEAMFS v1 Theorem IV.1 (2026-04-28)}\label{ssec:falsification}

\paragraph{Setup.}
A BEAMFS v1 reference implementation~\cite{desbrieres2026beamfsv1}
was instrumented with RadFI v0.1.0 in a Yocto-based qemu-aarch64 lab
environment (Linux kernel 7.0.1, distro \texttt{poky-hardened} 1.0).
The filesystem was formatted with the v1 INODE\_UNIVERSAL protection
scheme (\(\mathrm{RS}(255,239)\) coverage on inode metadata,
allocation bitmap, and superblock; data blocks unprotected by the RS
layer). RadFI was configured with \texttt{hook\_blk = 1},
\texttt{target\_dev = 1792} (the test loop device),
\texttt{inject\_on\_read = 0} (write-only model, v0.1.0 semantics),
and \texttt{probability = 1000} (one flip per thousand intercepted
writes).

\paragraph{Result.}
A workload of 1024 sequential 4\,KiB writes followed by a re-mount
and a checksummed read produced 7 corrupted blocks at the user level
(SHA-256 mismatch versus the pre-write fixture). The BEAMFS v1
recovery operator \(\mathcal{G}\), invoked on read, did not return
any of these blocks to consistency, nor did it transition to
fail-closed: it returned the corrupted bytes silently to user space.
This is a counter-example
to~\cite[Theorem~IV.1]{desbrieres2026beamfsv1}: there exists
\(s \models\) and \(e \in \mathcal{E}_{\text{RadFI v0.1.0}}\) such
that \(\mathcal{G}^k(e(s)) \not\models\) and
\(\mathcal{G}^k(e(s)) \neq \bot\) for all \(k \in \mathbb{N}\).

\paragraph{Diagnosis.}
The v1 INODE\_UNIVERSAL scheme, by design, did not extend RS
coverage to data blocks. The recovery operator's coverage of
\(\mathcal{E}_{\text{SEE}}\)
in~\cite[Theorem~IV.1]{desbrieres2026beamfsv1} was conditional on
the layout's redundancy specification (invariant I1 of the v1
report), which the v1 layout did not satisfy on data blocks. The
RadFI experiment exposed this incompleteness experimentally rather
than theoretically.

\paragraph{Consequence.}
The BEAMFS author retracted the wording of v1 Theorem IV.1 on the
same date and revised the threat model from radiation-only (v1) to
full electromagnetic resilience (v2 EMR), introducing the INLINE
protection scheme (\(\mathrm{RS}(255,239) \times 16\) per-block on
data blocks unconditionally)~\cite{desbrieres2026beamfsv2}. The v2
successor report states and proves a revised Theorem v2.1 covering
data-block perturbations.

This falsification is the central empirical contribution of the
present report. It is consistent with the Popperian
principle~\cite{popper1959logic} that a soundness theorem stands
until a single counter-example is produced, and that the
counter-example must motivate, not merely refute, the next theorem.

\subsection{Confirmation of BEAMFS v2 INLINE Recovery (2026-04-29)}\label{ssec:confirmation}

\paragraph{Setup.}
A BEAMFS v2 INLINE-formatted filesystem (\(\mathrm{RS}(255,239)
\times 16\), per-block protection on all data blocks) was created
with \texttt{mkfs.beamfs -s inline} on a 64\,MiB loop image. A
conformance fixture canary block was emitted at logical block 19 by
the mkfs tool, with byte-deterministic content yielding SHA-256
\(=\)
\texttt{\seqsplit{ced446d9bf5e6682a48e8782ce5ad33f575f08921451afaa127f26dc37515bab}}
for the user-visible 3{,}824-byte payload, and SHA-256 \(=\)
\texttt{\seqsplit{fb8a3b9e2704ce3fc11a0d0a4139d36c916cb5e7230acde8d2fa00ace739a91c}}
for the raw 4{,}096-byte on-disk block, both reproducible across
x86\_64 and aarch64 builds.

\paragraph{Method.}
Direct disk-level corruption (not RadFI) was applied: a single byte
at offset 77{,}824 (logical block 19, header byte 0) was XOR-flipped
from \texttt{0x42} to \texttt{0x43}. The image was then mounted
through BEAMFS, page caches were dropped, and the canary file was
read through the standard VFS path.

\paragraph{Result.}
The user-level read returned the canonical 3{,}824-byte payload with
SHA-256 matching the fixture. The kernel log recorded
\texttt{beamfs/inline: ino=2 iblock=0 subblock=0: 1 symbol(s) corrected},
confirming the RS decoder corrected one symbol error on subblock 0
of the canary block. A subsequent raw read of the on-disk block,
after a second cache drop, returned the canonical 4{,}096-byte
content with SHA-256 matching the fixture: the autonomic in-place
repair had persisted the correction byte-perfect to disk. The
recovery operator's idempotence (BEAMFS invariant I4 of the v2 report~\cite{desbrieres2026beamfsv2}) was confirmed on a third read
producing identical results without further log entries.

\paragraph{Methodological note.}
This experiment used direct disk-level corruption rather than RadFI
bio-layer injection. The choice was operationally driven: during
the same session, RadFI v0.1.2 \texttt{hook\_blk} pre-handler was
observed to register on \texttt{submit\_bio\_noacct} (kprobe armed
at the correct address, native kprobe events registering 10 hits
per \texttt{cat} invocation) but the RadFI counters did not
increment. Diagnosis is deferred to v2 of this report. The two
methods are equivalent at the SEU semantic level: each produces a
single bit-flip on volatile or persistent memory in the read path.
Direct disk corruption was the simpler chain of trust for the v2
INLINE empirical result, and the result is paper-grade reproducible
(fixture SHA-256 values are byte-deterministic across platforms).

\subsection{Falsification Matrix}

\begin{table*}[t]
\caption{Empirical falsification matrix. The v1 result is a single
counter-example sufficient to retract Theorem IV.1
of~\cite{desbrieres2026beamfsv1}. The v2 result is a single positive
instance corroborating, but not proving, the revised v2 recovery
properties.}
\label{tab:falsification-matrix}
\centering
\small
\begin{tabular}{@{}p{3cm}p{4.2cm}p{3.8cm}p{5.3cm}@{}}
\toprule
Implementation & Protection Scheme & Perturbation & Result \\
\midrule
BEAMFS v1 (2026-04-28) & INODE\_UNIVERSAL (metadata only) & RadFI v0.1.0, write hook, 1000\,PPM & \textbf{7/1024 silent corruption} (falsification of v1 Theorem IV.1) \\
BEAMFS v2 (2026-04-29) & UNIVERSAL\_INLINE (RS\(\times\)16 per data block) & \texttt{dd}-based single byte flip on canary block & \textbf{byte-perfect recovery + autonomic repair} (corroboration of pending v2 Theorem) \\
\bottomrule
\end{tabular}
\end{table*}
